% !TeX root = RJwrapper.tex
\title{visStatistics: Automated selection and visualisation of statistical hypothesis tests in R}


\author{by Sabine Schilling}

\maketitle

\abstract{%
We present \CRANpkg{visStatistics}, an R package that automates the selection and visualisation of statistical hypothesis tests between two vectors. Given vectors of class \code{numeric}, \code{integer}, or \code{factor}, the main function \code{visstat()} determines the appropriate test from the data types, distributional assumptions, and sample size alone, without any manual intervention. Visual outputs --- box plots, bar charts, regression lines with confidence bands, mosaic plots, and residual diagnostics --- are annotated with test statistics, assumption-check results, and post-hoc comparisons where applicable. The decision logic follows the general linear model framework, assessing normality on pooled model residuals rather than per group, which preserves statistical power. The package is particularly suited for server-side R applications where end users interact exclusively via a web interface, and for time-constrained settings such as statistical consulting or introductory teaching. All tests are applied using their default settings from base R, keeping the dependency footprint minimal and the decision tree transparent.
}

\section{Introduction}\label{introduction}

Most routine data analyses reduce to a comparatively small set of inferential
frameworks: group comparisons, contingency-table analyses, correlation, and
simple regression \citep{Sato:2017, Chicco:2025}. Selecting the appropriate test
within each framework, however, requires evaluating variable types, sample
sizes, and distributional assumptions --- a process that is error-prone,
time-consuming, and a recurring bottleneck in statistical consulting and
undergraduate teaching.

Several R packages address parts of this problem.
\CRANpkg{compareGroups} \citep{Subirana:2014} automates test choice via a
per-group Shapiro--Wilk normality test and produces publication-ready
descriptive tables, but its primary output is tabular rather than graphical,
and assumption diagnostics are not displayed alongside results.
\CRANpkg{ggstatsplot} \citep{Patil:2021} produces information-rich
\CRANpkg{ggplot2}-based \citep{Wickham:2016} plots with embedded test results, but
requires the user to select the appropriate plotting function for each
data-type combination (e.g., \texttt{ggbetweenstats()} for group comparisons,
\texttt{ggscatterstats()} for regression). \CRANpkg{rstatix} \citep{Kassambara:2025}
provides a pipe-friendly framework for common tests with tidy output, and
\CRANpkg{ggpubr} \citep{Kassambara:2026} adds statistical annotations to
publication-ready plots; both leave test selection entirely to the user.

\CRANpkg{visStatistics} takes a different approach: a single function
\texttt{visstat(x,\ y)} accepts any combination of numeric and categorical vectors
and resolves the full decision pipeline --- test selection, execution,
assumption checking, and visualisation --- without any further input. Three
design choices distinguish it from the packages above.

First, normality is assessed on the pooled residuals of a linear model
\texttt{lm(y\ \textasciitilde{}\ x)}, following the general linear model (GLM) framework
\citep{Searle:1971}, rather than per group. Group-wise testing is statistically
less powerful and more likely to reject normality unnecessarily, pushing the
analysis to non-parametric methods \citep{Lumley:2002}. When heteroscedasticity is
subsequently detected, group-wise normality diagnostics are shown additionally,
since Welch's methods use group-specific variance estimates.

Second, assumption diagnostics are always displayed, regardless of whether
assumptions are met. This enables users to visually override the automated
decision in individual cases, and supports the pedagogical goal of making
assumption checking an integral part of the analysis rather than an
afterthought.

Third, all tests use their default settings from base R functions, keeping
the dependency footprint minimal and the decision tree transparent. The
package does not wrap external statistical engines; the automated layer is
purely in the routing logic.

\section{Package overview}\label{package-overview}

\CRANpkg{visStatistics} is available on CRAN. The package accepts input in
three equivalent forms:

\begin{verbatim}
# Recommended form:
visstat(x, y)

# Formula interface:
visstat(y ~ x, data = dataframe)

# Backward-compatible form:
visstat(dataframe, "namey", "namex")
\end{verbatim}

where \texttt{x} and \texttt{y} are vectors of class \texttt{"numeric"}, \texttt{"integer"}, or
\texttt{"factor"}. The function returns an object of class \texttt{"visstat"} with
\texttt{print()}, \texttt{summary()}, and \texttt{plot()} S3 methods, and optionally writes
graphics to disk via \CRANpkg{Cairo} \citep{Urbanek:2025}.

\section{Decision logic}\label{decision-logic}

The significance level \(\alpha = 1 - \texttt{conf.level}\) is user-controlled
(default 0.95). Figure \ref{fig:overview} gives a high-level overview of all
implemented tests by input class; Figure \ref{fig:decision-tree} details the
decision path for a numeric response and categorical predictor. Technical
definitions of all test statistics are collected in \hyperref[tests]{Appendix C}.

\begin{figure}
\includegraphics[width=1\linewidth,alt={Flowchart showing all implemented statistical tests organised by the class of the input vectors.}]{figures/overview} \caption{Overview of all implemented tests by input class.}\label{fig:overview}
\end{figure}

\begin{figure}
\includegraphics[width=1\linewidth,alt={Decision tree selecting among Welch t-test, Student t-test, Wilcoxon, Fisher ANOVA, Welch ANOVA, and Kruskal-Wallis tests based on sample size, Shapiro-Wilk, and Levene test results.}]{figures/decision_tree} \caption{Decision tree for a numeric response and categorical predictor, based on sample size, normality of residuals, and homoscedasticity.}\label{fig:decision-tree}
\end{figure}

\subsection{Numeric response and categorical predictor: comparing central tendencies}\label{numeric-response-and-categorical-predictor-comparing-central-tendencies}

A linear model \texttt{lm(y\ \textasciitilde{}\ x)} is fitted and assumption diagnostics are always
displayed via \texttt{vis\_lm\_assumptions()} (see \hyperref[glm]{Appendix A}).
The Shapiro--Wilk test \citep{Shapiro:1965} is applied to internally studentised
residuals from \texttt{rstandard()} \citep{Cook:1982}. When all groups contain more than
50 observations, normality is assumed by the central limit theorem
\citep{Lumley:2002, Rasch:2011}.

If normality is rejected (\(p_\text{SW} \le \alpha\)), non-parametric tests are
selected: \texttt{wilcox.test()} \citep{Mann:1947} for two groups, or \texttt{kruskal.test()}
\citep{Kruskal:1952} followed by Holm-adjusted \texttt{pairwise.wilcox.test()} \citep{Holm:1979}
for more than two groups.

If normality is not rejected, the Levene--Brown--Forsythe test \texttt{levene.test()}
\citep{Brown:1974} (see \hyperref[variance]{Appendix B}) assesses variance homogeneity. For
homoscedastic data, Student's \texttt{t.test(var.equal\ =\ TRUE)} or Fisher's \texttt{aov()}
\citep{Fisher:1990} with \texttt{TukeyHSD()} \citep{Hochberg:1987} is applied. For
heteroscedastic data, Welch's \texttt{t.test()} \citep{Welch:1947, Satterthwaite:1946}
or \texttt{oneway.test()} \citep{Welch:1951} with \texttt{games.howell()} \citep{Games:1976} is
applied. When switching to Welch methods, group-wise normality diagnostics are
additionally displayed via \texttt{vis\_group\_normality()}, since the pooled \texttt{lm()}
residuals are no longer the appropriate diagnostic under heteroscedasticity.

\subsection{Both variables numeric: regression or rank correlation}\label{both-variables-numeric-regression-or-rank-correlation}

By default, \texttt{visstat()} fits a simple linear regression via \texttt{lm()}, regardless
of whether GLM assumptions are met, and always displays assumption diagnostics.
When \texttt{correlation\ =\ TRUE} is set explicitly, Spearman's \(\rho\) is computed via
\texttt{cor.test(...,\ method\ =\ "spearman")}. Because the choice between modelling a
directional relationship and measuring a monotone association cannot be
automated from data types alone, Spearman correlation is never triggered
automatically.

\subsection{Both variables categorical: comparing proportions}\label{both-variables-categorical-comparing-proportions}

When both variables are categorical, \texttt{visstat()} tests independence using
Pearson's \(\chi^2\) test (\texttt{chisq.test()}) or Fisher's exact test
(\texttt{fisher.test()}), depending on expected cell counts following Cochran's rule
\citep{Cochran:1954}: the \(\chi^2\) approximation is used if no expected cell count
is less than 1 and no more than 20\% of cells have expected counts below 5.
Yates' continuity correction \citep{Yates:1934} is applied by default to
\(2 \times 2\) tables. For general \(R \times C\) tables, Pearson's \(\chi^2\) result
is supplemented with a mosaic plot from \CRANpkg{vcd} \citep{Meyer:2006, Meyer:2024}
with tiles coloured by Pearson residuals.

\subsection{Ordered factors}\label{ordered-factors}

When the response is an ordered factor, \texttt{visstat()} converts it to integer
level codes and follows the non-parametric path. When both variables are
ordered and \texttt{correlation\ =\ TRUE}, Kendall's \(\tau_b\) \citep{Kendall:1945, Agresti:2010}
is computed instead, as it corrects for ties explicitly --- unavoidable with
few ordinal levels \citep{Xu:2013}.

\section{Examples}\label{examples}

The number of plots returned by \texttt{visstat()} depends on the input types and
the outcome of assumption checks. Tests that have no separate assumption
diagnostics --- Kendall's \(\tau_b\) for two ordered factors --- produce a single
plot. Most tests produce two panels: an assumption-diagnostic panel and a
result panel. When heteroscedasticity is detected and \texttt{visstat()} routes to a
Welch method, a third panel with group-wise normality diagnostics is added,
since pooled-residual diagnostics are no longer appropriate under unequal
variances.

\subsection{Two groups: Welch's t-test}\label{two-groups-welchs-t-test}

With binary factor \texttt{am} (transmission type, 0 = automatic, 1 = manual) and
continuous response \texttt{mpg} (fuel efficiency) from the \texttt{mtcars} dataset,
\texttt{visstat()} fits \texttt{lm(mpg\ \textasciitilde{}\ am)}, applies Shapiro--Wilk to the pooled residuals
(not rejected), and subsequently the Levene--Brown--Forsythe test, which detects
heteroscedasticity. The routing therefore leads to Welch's t-test rather than
Student's t-test. Because the pooled-residual diagnostic is no longer
appropriate under heteroscedasticity, \texttt{vis\_group\_normality()} additionally
displays per-group normality checks, yielding three panels in total.

\begin{verbatim}
mtcars$am <- as.factor(mtcars$am)
t_test_stats <- visstat(mtcars$am, mtcars$mpg)
\end{verbatim}

\begin{figure}
\includegraphics[width=0.32\linewidth]{visStatistics-RJournal_files/figure-latex/ttest-example-1} \includegraphics[width=0.32\linewidth]{visStatistics-RJournal_files/figure-latex/ttest-example-2} \includegraphics[width=0.32\linewidth]{visStatistics-RJournal_files/figure-latex/ttest-example-3} \caption{Welch's t-test applied to the \texttt{mtcars} dataset (\texttt{am} vs.\ \texttt{mpg}). Pooled-residual assumption diagnostics (Shapiro--Wilk, Bartlett, Levene), group-wise normality diagnostics, and box plots of fuel efficiency by transmission type with the Welch t-test result.}\label{fig:ttest-example}
\end{figure}

\subsection{Multiple groups: one-way ANOVA with Tukey HSD}\label{multiple-groups-one-way-anova-with-tukey-hsd}

With factor \texttt{block} (experimental block) and continuous response \texttt{yield} from
the \texttt{npk} pea-growing dataset, Shapiro--Wilk does not reject normality of
pooled residuals and the Levene test does not reject homoscedasticity. Fisher's
one-way ANOVA is therefore applied, followed by Tukey HSD post-hoc comparisons
with a compact letter display.

\begin{verbatim}
anova_npk <- visstat(npk$block, npk$yield, conf.level = 0.95)
\end{verbatim}

\begin{figure}
\includegraphics[width=0.48\linewidth]{visStatistics-RJournal_files/figure-latex/anova-example-1} \includegraphics[width=0.48\linewidth]{visStatistics-RJournal_files/figure-latex/anova-example-2} \caption{Fisher's one-way ANOVA applied to the \texttt{npk} dataset (\texttt{block} vs.\ \texttt{yield}). Assumption diagnostics and box plots of pea yield by experimental block; compact letter display (TukeyHSD, $\alpha = 0.05$) shows no significant pairwise differences.}\label{fig:anova-example}
\end{figure}

\subsection{Numeric response and numeric predictor: simple linear regression}\label{numeric-response-and-numeric-predictor-simple-linear-regression}

With trunk \texttt{Girth} and \texttt{Volume} from the \texttt{trees} dataset, both vectors are
numeric, so \texttt{visstat()} fits \texttt{lm(Volume\ \textasciitilde{}\ Girth)} and displays confidence and
prediction bands at the user-specified level \texttt{conf.level\ =\ 0.9}. Assumption
diagnostics include the Shapiro--Wilk and Anderson--Darling tests for residual
normality and the Breusch--Pagan test for homoscedasticity (see
\hyperref[variance]{Appendix B}).

\begin{verbatim}
linreg_trees <- visstat(trees$Girth, trees$Volume, conf.level = 0.9)
\end{verbatim}

\begin{verbatim}
#> Warning: Statistical assumptions violated:
#> Homoscedasticity violated (Breusch-Pagan p = 0.0178 )
#> Analysis proceeded but interpret results cautiously.
\end{verbatim}

\begin{verbatim}
#> RECOMMENDATION: Consider exploring alternatives outside visstat() such as data transformations, generalised linear models, or robust regression. For an assumption-free, non-causal alternative consider rerunning with correlation = TRUE.
\end{verbatim}

\begin{figure}
\includegraphics[width=0.48\linewidth]{visStatistics-RJournal_files/figure-latex/regression-example-1} \includegraphics[width=0.48\linewidth]{visStatistics-RJournal_files/figure-latex/regression-example-2} \caption{Simple linear regression of \texttt{Volume} on \texttt{Girth} for the \texttt{trees} dataset (\texttt{conf.level = 0.9}). Assumption diagnostics (Shapiro--Wilk, Anderson--Darling, Breusch--Pagan) and scatter plot with fitted regression line, 90\% confidence band (dark shading), and 90\% prediction band (light shading).}\label{fig:regression-example}
\end{figure}

\subsection{\texorpdfstring{Two categorical variables: Pearson's \(\chi^2\) test}{Two categorical variables: Pearson's \textbackslash chi\^{}2 test}}\label{two-categorical-variables-pearsons-chi2-test}

With \texttt{Eye} and \texttt{Hair} from the \texttt{HairEyeColor} dataset, all expected cell
counts exceed the Cochran thresholds \citep{Cochran:1954}, so the \(\chi^2\)
approximation is used. For this \(4 \times 4\) table, \texttt{visstat()} returns both a
bar chart of row percentages and a mosaic plot with tiles coloured by Pearson
residuals, making the cells driving the association immediately visible.

\begin{verbatim}
hair_eye_df <- counts_to_cases(as.data.frame(HairEyeColor))
visstat(hair_eye_df$Eye, hair_eye_df$Hair)
\end{verbatim}

\begin{figure}
\includegraphics[width=0.48\linewidth]{visStatistics-RJournal_files/figure-latex/chisq-example-1} \includegraphics[width=0.48\linewidth]{visStatistics-RJournal_files/figure-latex/chisq-example-2} \caption{Pearson's $\chi^2$ test applied to the \texttt{HairEyeColor} dataset. Grouped bar chart of eye colour by hair colour and mosaic plot with tiles coloured by Pearson residuals (blue: over-represented, red: under-represented).}\label{fig:chisq-example}
\end{figure}

\section{Saving graphical output}\label{saving-graphical-output}

All plots can be saved in any format supported by \CRANpkg{Cairo}
\citep{Urbanek:2025} --- \texttt{"png"}, \texttt{"pdf"}, \texttt{"svg"}, \texttt{"ps"}, \texttt{"tiff"} --- by supplying
the \texttt{graphicsoutput} and \texttt{plotDirectory} arguments. File names follow the
pattern \texttt{testname\_namey\_namex.ext} by default, or a user-supplied
\texttt{plotName}. The full file paths of generated graphics are stored as the
\texttt{"plot\_paths"} attribute on the returned \texttt{"visstat"} object.

\section{Limitations}\label{limitations}

\subsection{Scope of automated test selection}\label{scope-of-automated-test-selection}

All tests are applied with their default settings from base R. As a
consequence, paired tests are not supported and interaction terms are not
available in the ANOVA path. Only simple linear regression is implemented, as
multiple regression cannot be visualised in two dimensions. When regression
assumptions are violated, \CRANpkg{visStatistics} offers Spearman rank
correlation (\texttt{correlation\ =\ TRUE}) as the sole automated alternative.
Alternative methods such as data transformations, generalised linear models, or
robust regression are not implemented: each requires user judgment --- about the
transformation family, the link function, or the estimator --- that cannot be
automated without substantially expanding the decision tree and increasing the
risk of Type I error inflation.

\subsection{P-value-based assumption testing}\label{p-value-based-assumption-testing}

No single test maintains optimal Type I error rates and power across all
distributions \citep{Olejnik:1987}. Normality tests behave poorly at both extremes
of the sample-size range: they fail to detect non-normality in small samples
and flag negligible departures as significant in large ones \citep{Ghasemi:2012, Fagerland:2012}. The \(n > 50\) per-group threshold for bypassing the
Shapiro--Wilk test is necessarily arbitrary; simulation studies indicate that
the required sample size for the central limit theorem to take effect depends
on the skewness and kurtosis of the underlying distribution \citep{Fagerland:2012}.
Assumption tests provide no information on the nature of the deviation
\citep{Shatz:2024}, and combining multiple tests inflates the overall Type I error
rate. For these reasons, the diagnostic plots generated by \texttt{visstat()} should
always be inspected alongside the automated decision; the \texttt{conf.level}
argument provides a direct handle on the \(\alpha\) used throughout.

\subsection{Visualisation}\label{visualisation}

The package depends exclusively on base R graphics with no
\CRANpkg{ggplot2} \citep{Wickham:2016} dependency, keeping the transitive
dependency footprint minimal. For more polished, annotated plots of chosen
statistical tests, see \CRANpkg{ggstatsplot} \citep{Patil:2021} or
\CRANpkg{ggpubr} \citep{Kassambara:2026}.

\section{Summary}\label{summary}

\CRANpkg{visStatistics} provides a single-function interface to automated
statistical test selection, execution, and visualisation for the most
commonly used inferential frameworks in two-variable analyses. Its principal
contribution is the integration of the full pipeline --- data-type detection,
distributional assumption checking on GLM residuals, test routing, and
annotated graphical output --- into a reproducible, server-deployable workflow
with a minimal dependency footprint. The package is available on CRAN at
\texttt{https://CRAN.R-project.org/package=visStatistics} and on GitHub at
\texttt{https://github.com/shhschilling/visStatistics}.

\section{Appendix A: The general linear model}\label{glm}

The general linear model \citep{Searle:1971} provides a unified framework
underlying Student's t-test, Fisher's ANOVA, and simple linear regression:

\[Y_i = \beta_0 + \beta_1 x_{i1} + \cdots + \beta_{k-1} x_{i,k-1} + \varepsilon_i,
\quad \varepsilon_i \sim N(0, \sigma^2).\]

\textbf{Student's t-test} uses one binary indicator variable \(x_{i1}\); testing
\(H_0: \beta_1 = 0\) is equivalent to testing \(H_0: \mu_1 = \mu_2\).

\textbf{Fisher's ANOVA} uses \(k-1\) indicator variables for \(k\) groups; testing
\(H_0: \beta_1 = \cdots = \beta_{k-1} = 0\) is equivalent to testing equality
of all group means.

\textbf{Simple linear regression} uses one continuous predictor; testing
\(H_0: \beta_1 = 0\) examines whether a linear relationship exists.

The two-sample case connects t-tests and ANOVA: \(t^2 = F\). For two groups,
\texttt{t.test(var.equal\ =\ TRUE)} and \texttt{aov()} yield identical \(p\)-values, as do
\texttt{t.test(var.equal\ =\ FALSE)} and \texttt{oneway.test()}.

\section{Appendix B: Tests for homoscedasticity}\label{variance}

\subsection{\texorpdfstring{The Levene--Brown--Forsythe test \texttt{levene.test()}}{The Levene--Brown--Forsythe test levene.test()}}\label{the-levenebrownforsythe-test-levene.test}

The test \citep{Brown:1974} evaluates equality of variances using the absolute
deviations from group medians \(z_{ij} = |y_{ij} - \tilde{y}_i|\) and applies
a one-way ANOVA \(F\)-statistic to the \(z_{ij}\). Using the median rather than
the mean makes it robust to skewed data \citep{Allingham:2012}. The implementation
\texttt{levene.test()} mimics the default behaviour of \texttt{leveneTest()} in
\CRANpkg{car} \citep{Fox:2019} and drives the parametric/Welch branching decision.

\subsection{\texorpdfstring{Bartlett's test \texttt{bartlett.test()}}{Bartlett's test bartlett.test()}}\label{bartletts-test-bartlett.test}

Bartlett's test \citep{Bartlett:1937} has greater power than the Brown--Forsythe
version when normality holds \citep{Allingham:2012}. Its statistic

\[K^2 = \frac{(N-k)\ln s_p^2 - \sum_{i=1}^k (n_i-1)\ln s_i^2}
{1 + \frac{1}{3(k-1)}\!\left(\sum_{i=1}^k \frac{1}{n_i-1} - \frac{1}{N-k}\right)}\]

follows \(\chi^2(k-1)\) under the null. Bartlett's test result is shown in the
assumption plot but does \textbf{not} drive the branching decision; only
\texttt{levene.test()} does.

\subsection{\texorpdfstring{Breusch--Pagan test \texttt{bp.test()}}{Breusch--Pagan test bp.test()}}\label{breuschpagan-test-bp.test}

For simple linear regression, group-based variance tests are not applicable.
The Breusch--Pagan test \citep{Breusch:1979} assesses whether the variance of the
residuals depends on the predictor values and is shown in the regression
assumption plot.

\section{Appendix C: Test statistics}\label{tests}

\subsection{Student's t-test and Welch's t-test}\label{students-t-test-and-welchs-t-test}

The test statistic for Student's t-test is

\[t = \frac{\bar{x}_1 - \bar{x}_2}{s_p \sqrt{\frac{1}{n_1} + \frac{1}{n_2}}},
\quad s_p = \sqrt{\frac{(n_1-1)s_1^2 + (n_2-1)s_2^2}{n_1+n_2-2}},\]

following a \(t\)-distribution with \(n_1 + n_2 - 2\) degrees of freedom.
Welch's t-test relaxes the homoscedasticity assumption:

\[t = \frac{\bar{x}_1 - \bar{x}_2}{\sqrt{s_1^2/n_1 + s_2^2/n_2}},\]

with degrees of freedom approximated by the Welch--Satterthwaite equation
\citep{Welch:1947, Satterthwaite:1946}. Welch's methods outperform their classical
counterparts when variances differ \citep{Moser:1992, Fagerland:2009, Delacre:2017}.

\subsection{Wilcoxon rank-sum test}\label{wilcoxon-rank-sum-test}

The test statistic is the Mann--Whitney \(U\) statistic \citep{Mann:1947}:

\[W = U_1 = W_1 - \frac{n_1(n_1+1)}{2},\]

where \(W_1 = \sum_{i=1}^{n_1} R(x_{1,i})\) is the rank sum of the first group.
An exact \(p\)-value is computed for samples of fewer than 50 observations
without ties; otherwise a normal approximation with continuity correction is
used.

\subsection{Fisher's one-way ANOVA and Welch's heteroscedastic ANOVA}\label{fishers-one-way-anova-and-welchs-heteroscedastic-anova}

Fisher's ANOVA test statistic is

\[F = \frac{SS_\text{between}/(k-1)}{SS_\text{within}/(N-k)},\]

following an \(F(k-1, N-k)\) distribution. Post-hoc comparisons use
\texttt{TukeyHSD()} \citep{Hochberg:1987}, which controls the family-wise error rate via
the studentised range distribution \citep{Abdi:2007}.

Welch's heteroscedastic ANOVA (\texttt{oneway.test()}) \citep{Welch:1951} uses
group-specific variance weights and a Satterthwaite-type adjustment for degrees
of freedom. Post-hoc comparisons use \texttt{games.howell()} \citep{Games:1976}, which
applies separate variance estimates per pair. Compact letter displays are
produced by \texttt{multcompLetters()} from \CRANpkg{multcompView} \citep{Graves:2026}.

\subsection{Kruskal--Wallis test}\label{kruskalwallis-test}

When normality is rejected, \texttt{kruskal.test()} \citep{Kruskal:1952} tests whether
the groups come from the same population based on ranks:

\[H = \frac{12}{N(N+1)} \sum_{i=1}^{k} n_i \left(\bar{R}_i - \bar{R}\right)^2,\]

approximately \(\chi^2(k-1)\) under the null. Post-hoc comparisons use
Holm-adjusted \texttt{pairwise.wilcox.test()} \citep{Holm:1979}.

\subsection{Simple linear regression and Spearman rank correlation}\label{simple-linear-regression-and-spearman-rank-correlation}

The regression plot displays the fitted line \(\hat{y} = b_0 + b_1 x\) with
pointwise confidence and prediction bands at the specified \texttt{conf.level}.
Assumption diagnostics always include the Shapiro--Wilk and Anderson--Darling
\citep{Gross:2015} tests for residual normality and the Breusch--Pagan test
\texttt{bp.test()} \citep{Breusch:1979} for homoscedasticity (see \hyperref[variance]{Appendix B}).

When \texttt{correlation\ =\ TRUE}, Spearman's \(\rho\) is computed as Pearson's \(r\)
applied to the ranks:

\[\rho = r\!\left(\operatorname{rank}(x),\, \operatorname{rank}(y)\right).\]

\subsection{\texorpdfstring{Pearson's \(\chi^2\) test and Fisher's exact test}{Pearson's \textbackslash chi\^{}2 test and Fisher's exact test}}\label{pearsons-chi2-test-and-fishers-exact-test}

Let \(O_{ij}\) and \(E_{ij}\) denote the observed and expected frequencies in row
\(i\) and column \(j\) of an \(R \times C\) contingency table. The Pearson residual
for cell \((i,j)\) is

\[r_{ij} = \frac{O_{ij} - E_{ij}}{\sqrt{E_{ij}}}.\]

The test statistic of Pearson's \(\chi^2\) test \citep{Pearson:1900} is

\[\chi^2 = \sum_{i,j} r_{ij}^2,\]

compared to \(\chi^2\!\left((R-1)(C-1)\right)\). For general \(R \times C\)
tables, \texttt{visstat()} supplements the bar chart with a mosaic plot in which
tiles are coloured by \(r_{ij}\) (blue: positive, red: negative), making the
cells driving the association immediately visible.

Fisher's exact test \citep{Fisher:1970} is applied when Cochran's rule
\citep{Cochran:1954} is violated. It evaluates the exact hypergeometric probability
of the observed table given the fixed margins.

\subsection{\texorpdfstring{Kendall's \(\tau_b\)}{Kendall's \textbackslash tau\_b}}\label{kendalls-tau_b}

For two ordinal variables with \(n\) joint observations, Kendall's \(\tau_b\)
\citep{Kendall:1945} is

\[\tau_b = \frac{C - D}{\sqrt{(n_0 - n_1)(n_0 - n_2)}},\]

where \(C\) and \(D\) are the numbers of concordant and discordant pairs,
\(n_0 = n(n-1)/2\), and \(n_1\), \(n_2\) correct for ties in the response and
predictor respectively. Kendall's \(\tau_b\) is preferred over Spearman's \(\rho\)
for ordinal data with few levels, where ties are common \citep{Agresti:2010, Xu:2013}.

\bibliography{Rjreferences.bib}

\address{%
Sabine Schilling\\
Institute of Tourism and Mobility, Lucerne University of Applied Sciences and Arts\\%
Rösslimatte 48, 6002 Lucerne, Switzerland\\
%
\url{https://github.com/shhschilling/visStatistics}\\%
\textit{ORCiD: \href{https://orcid.org/0000-0002-8318-9421}{0000-0002-8318-9421}}\\%
\href{mailto:sabineschilling@gmx.ch}{\nolinkurl{sabineschilling@gmx.ch}}%
}
